%% ====================================================================
\section{Evaluation}
\label{sec:evaluation}
%% ====================================================================

We evaluate the \MLIPsOntology{} along multiple dimensions: competency
question coverage via SPARQL queries, consistency verification via OWL
reasoning, automated pitfall detection, and comparison with existing
ontologies.

\subsection{Competency Question Coverage}
\label{sec:eval-cq}

We demonstrate that each competency question (Section~\ref{sec:requirements})
is answerable as a SPARQL query over the ontology. For each CQ, we provide the
query template; Appendix~\ref{sec:appendix-dl-queries} additionally presents
the equivalent description logic formulations.

%% TODO: populate with actual query results once the KG has data

\paragraph{CQ1: Algorithm hyperparameters.}
Which MLIP algorithms exist, and what hyperparameters does each accept?
\begin{lstlisting}[language=sparql]
SELECT ?algo ?hp ?name ?type ?default WHERE {
  ?algo a mlips:MLIPAlgorithm ;
        mlips:hasHyperparameter ?hp .
  ?hp mlips:hyperparameterName ?name ;
      mlips:hyperparameterDatatype ?type .
  OPTIONAL { ?hp mlips:defaultValue ?default }
}
\end{lstlisting}

\paragraph{CQ2: Implementations and versions.}
Which libraries implement a given algorithm?
\begin{lstlisting}[language=sparql]
SELECT ?algo ?impl ?lib ?version WHERE {
  ?algo a mlips:MLIPAlgorithm ;
        mlips:hasImplementation ?impl .
  ?impl mlips:implementedIn ?lib ;
        mlips:version ?version .
}
\end{lstlisting}

\paragraph{CQ3: Training datasets for a material system.}
For a given material, which training datasets are available and what DFT
settings were used?
\begin{lstlisting}[language=sparql]
SELECT ?ds ?xc ?cutoff ?pseudo WHERE {
  ?ds a mlips:TrainingDataset ;
      mlips:coversMaterial ?mat .
  ?mat mlips:chemicalFormula "TiAl" .
  ?ds mlips:hasDFTCalculation ?calc .
  ?calc mlips:hasDFTSettings ?settings .
  ?settings mlips:xcFunctional ?xc ;
            mlips:energyCutoff ?cutoff ;
            mlips:pseudopotentialType ?pseudo .
}
\end{lstlisting}

\paragraph{CQ4: Dataset provenance.}
What is the provenance of a training dataset?
\begin{lstlisting}[language=sparql]
SELECT ?ds ?label ?prov WHERE {
  ?ds a mlips:TrainingDataset ;
      rdfs:label ?label ;
      mlips:datasetProvenance ?prov .
}
\end{lstlisting}

\paragraph{CQ5: Dataset size and property coverage.}
How many configurations does a dataset contain, and what properties are
covered?
\begin{lstlisting}[language=sparql]
SELECT ?ds ?nconfig ?prop WHERE {
  ?ds a mlips:TrainingDataset ;
      mlips:numConfigurations ?nconfig ;
      mlips:coversProperty ?prop .
}
\end{lstlisting}

\paragraph{CQ6: Published benchmarks for an algorithm and material.}
Which studies have benchmarked a given algorithm on a given material?
\begin{lstlisting}[language=sparql]
SELECT ?study ?result ?metric_type ?prop ?value WHERE {
  ?study a mlips:BenchmarkStudy ;
         mlips:hasResult ?result .
  ?result mlips:evaluatesModel ?model ;
          mlips:targetMaterial ?mat ;
          mlips:hasAccuracyMetric ?metric .
  ?model mlips:trainedWith ?algo .
  ?algo rdfs:label "MACE" .
  ?mat mlips:chemicalFormula "TiAl" .
  ?metric mlips:metricType ?metric_type ;
          mlips:metricProperty ?prop ;
          mlips:metricValue ?value .
}
\end{lstlisting}

\paragraph{CQ7: Best-performing algorithm--hyperparameter combinations.}
For a given material and target property, which combinations perform well?
\begin{lstlisting}[language=sparql]
SELECT ?algo ?value WHERE {
  ?result mlips:evaluatesModel ?model ;
          mlips:targetMaterial ?mat ;
          mlips:hasAccuracyMetric ?metric .
  ?model mlips:trainedWith ?algo .
  ?mat mlips:chemicalFormula "TiAl" .
  ?metric mlips:metricType mlips:RMSE ;
          mlips:metricProperty mlips:EnergyProperty ;
          mlips:metricValue ?value .
} ORDER BY ?value LIMIT 10
\end{lstlisting}

\paragraph{CQ8: Simulation types supported by an algorithm.}
Which simulation types have been performed with a given MLIP?
\begin{lstlisting}[language=sparql]
SELECT ?algo ?sim WHERE {
  ?algo a mlips:MLIPAlgorithm ;
        mlips:supportsSimulation ?sim .
}
\end{lstlisting}

\subsection{Consistency and Classification}
\label{sec:eval-consistency}

We verify the ontology's logical consistency using the HermiT OWL~2
reasoner~\cite{glimm2014hermit} via the ROBOT
tool~\cite{jackson2019robot}:

\begin{lstlisting}[language=bash,caption={Consistency verification.}]
robot reason --reasoner HermiT -i mlips.ttl
\end{lstlisting}

The reasoner confirms that:
\begin{itemize}
\item The ontology is \emph{consistent}: no contradictions exist among the
  axioms (A1)--(A27).
\item No class is \emph{unsatisfiable}: every class can have instances.
\item The inferred class hierarchy contains no unexpected subsumptions.
\end{itemize}

\subsection{Automated Pitfall Detection}
\label{sec:eval-oops}

We evaluate the ontology using OOPS! (OntOlogy Pitfall
Scanner)~\cite{poveda2014oops}, which checks for common ontology design
pitfalls. OOPS! classifies pitfalls into three severity levels: critical,
important, and minor.

%% TODO: insert OOPS! results once the web service is available
%% The OOPS! API at https://oops.linkeddata.es/rest is currently
%% returning server errors. Run manually via the web interface.

\DanielHernandez{Run OOPS! when the web service is back and insert results here.}

\subsection{Comparison with Existing Ontologies}
\label{sec:eval-comparison}

Table~\ref{tab:comparison} compares the \MLIPsOntology{} with existing
ontologies from the ML and materials science domains. For each competency
question, we indicate whether it is answerable ($\bullet$), partially
answerable ($\circ$), or not answerable (---) using each ontology.

\begin{table}[t]
\centering
\caption{Competency question coverage across ontologies. $\bullet$ =
  answerable, $\circ$ = partially, --- = not answerable.}
\label{tab:comparison}
\begin{tabular}{l@{\hspace{6pt}}c@{\hspace{6pt}}c@{\hspace{6pt}}c@{\hspace{6pt}}c@{\hspace{6pt}}c@{\hspace{6pt}}c}
\toprule
 & \rotatebox{70}{ML-Schema} & \rotatebox{70}{MDO} & \rotatebox{70}{CMSO/ASMO} & \rotatebox{70}{EMMO} & \rotatebox{70}{PMDco} & \rotatebox{70}{\textbf{MLIPs}} \\
\midrule
CQ1: Algo.\ hyperparams     & $\circ$ & ---     & ---     & ---     & ---     & $\bullet$ \\
CQ2: Implementations         & $\circ$ & ---     & ---     & ---     & ---     & $\bullet$ \\
CQ3: Training data + DFT     & ---     & $\circ$ & $\circ$ & $\circ$ & $\circ$ & $\bullet$ \\
CQ4: Dataset provenance       & ---     & ---     & ---     & ---     & ---     & $\bullet$ \\
CQ5: Dataset size + coverage  & ---     & ---     & ---     & ---     & ---     & $\bullet$ \\
CQ6: Published benchmarks     & ---     & ---     & ---     & ---     & ---     & $\bullet$ \\
CQ7: Best combinations        & ---     & ---     & ---     & ---     & ---     & $\bullet$ \\
CQ8: Simulation types         & ---     & ---     & $\circ$ & $\circ$ & ---     & $\bullet$ \\
\midrule
Training run provenance        & $\bullet$ & ---   & ---     & ---     & ---     & $\bullet$ \\
Crystal structure              & ---     & $\bullet$ & $\bullet$ & $\bullet$ & $\bullet$ & $\circ$ \\
\bottomrule
\end{tabular}
\end{table}

ML-Schema partially covers CQ1 and CQ2 through its general
$\mlsAlgorithmC$ and $\mlsImplementationC$ classes, but lacks
physics-specific hyperparameters and units. MDO and CMSO/ASMO partially
cover CQ3 through their calculation and structure modules, but do not model
MLIP-specific training data provenance or DFT settings at the level of
detail needed. No existing ontology covers CQ4--CQ7, which require the
benchmark--algorithm--dataset triangle that is central to the
\MLIPsOntology{}.

The last two rows show areas where existing ontologies complement ours:
ML-Schema's $\mlsRunC$ provides richer training run metadata (which we
align via $\TrainingRunC$), and CMSO provides detailed crystal structure
descriptions (which we link via $\MaterialSystemC$).

\subsection{Reusability}

The ontology is published under CC BY 4.0 with a persistent IRI and versioned
releases. SHACL shapes are distributed alongside the ontology. All resources
use standard Semantic Web formats (OWL/Turtle, SHACL, SPARQL), ensuring tool
independence.

\subsection{Adoption and Impact}

The resources are in active use within the META-LEARN project. The knowledge
graph is publicly queryable via a SPARQL endpoint. As the field grows, the
\MLIPsOntology{} provides the vocabulary for structured comparison and
discovery of MLIP studies---a capability not available from any existing
resource.

%% TODO: report concrete usage statistics

%% ====================================================================
